\documentclass[conference]{IEEEtran}
\IEEEoverridecommandlockouts

% --- packages -------------------------------------------------------------
\usepackage{amsmath,amssymb}
\usepackage{graphicx}
\usepackage{booktabs}
\usepackage{multirow}
\usepackage{siunitx}
\usepackage{xcolor}
\usepackage{microtype}
\usepackage[hidelinks]{hyperref}
\usepackage[capitalise,noabbrev]{cleveref}
\usepackage[textsize=scriptsize,color=orange!25,linecolor=orange!60]{todonotes}
\setlength{\marginparwidth}{2cm}

\graphicspath{{figures/}}
\sisetup{detect-weight=true, detect-family=true, group-minimum-digits=4}

% --- shorthands -----------------------------------------------------------
\newcommand{\Gprior}{G_{\mathrm{prior}}}
\newcommand{\Ghard}{G_{\mathrm{hardneg}}}
\newcommand{\Gtot}{G_{\mathrm{total}}}

\begin{document}

\title{Streaming Crossing-Onset Detection for Pedestrian\\Intention:
Protocol, Diagnosis, and Onset Timing Under Censoring}

\author{\IEEEauthorblockN{Author Name}
\IEEEauthorblockA{Affiliation \\ email@domain}}

\maketitle

\begin{abstract}
Pedestrian crossing-intention models are usually trained and tested under an
event-anchored protocol: each pedestrian track contributes one short observation
window, placed a known one to two seconds before an annotated crossing event. A
detector on a vehicle works differently. It watches continuously, and at every
moment it must judge whether a crossing is about to begin. This paper defines
both protocols precisely on the PIE dataset and measures what the difference
between them costs. Anchored sampling gives a class ratio of about $2.5{:}1$
because it takes roughly one window per track. Dense sliding-window sampling
gives $33.9{:}1$, because that is how often a crossing is genuinely imminent
when a camera watches without pause. The protocols also differ in what their
negatives are made of. We find that $25.7\%$ of streaming training windows and
$39.2\%$ of test windows show a pedestrian who does cross, only later than the
horizon asks about, and the anchored protocol produces no such windows for
either training or testing. A model trained under the anchored protocol falls
from $0.873$ to $0.540$ AUC when tested on the stream, and re-tuning its
decision threshold at the streaming base rate recovers almost none of that
loss. What the model has lost is therefore its ability to rank pedestrians by
risk, not the placement of its threshold. We argue that part of the difficulty
comes from the label itself, which gives opposite answers to almost identical
observations on either side of a fixed cutoff, and we propose treating the task
as predicting \emph{when} a crossing will start, under right-censoring. A
window whose future footage runs out is then a censored observation rather than
a negative, and the model still reports the probability of a crossing within the
original horizon, so it remains directly comparable with binary baselines.
\todo[inline]{Closing results sentence still missing: the onset formulation is
implemented and unit-tested but has not been trained. Do not submit without
\cref{tab:onset} filled in, or with the contribution restated as
formulation-only.}
\end{abstract}

\begin{IEEEkeywords}
pedestrian intention, crossing prediction, online action detection, streaming
evaluation, class imbalance, survival analysis
\end{IEEEkeywords}

% ==========================================================================
\section{Introduction}
\label{sec:intro}

A pedestrian-crossing detector on a moving vehicle is a streaming system. It
sees the road continuously, and at every moment it must judge whether someone is
about to step into it. Nothing tells it when to attend to a particular
pedestrian, and most of the time nothing is about to happen.

The protocol used to evaluate such systems asks a narrower question. Following
the procedure introduced with PIE and JAAD~\cite{rasouli2019pie,rasouli2017jaad}
and settled in the standard evaluation suite~\cite{kotseruba2021benchmark}, each
pedestrian track contributes about one labeled observation window. Observation
stops at the crossing event, and the window is placed so that its last frame
falls a known one to two seconds before that event. The class ratio that results
is mild, roughly $2.5{:}1$, and reported scores are high.

The two settings differ in more than difficulty. Because the anchored protocol
places its window relative to an event it already knows about, it never produces
the windows that occur earlier in a crossing pedestrian's track. Those are the
moments when a person who will eventually cross is still walking along the
sidewalk. Under a short horizon such windows are negatives, they look almost
exactly like the positives that follow them, and they are the situations in
which a deployed detector must not raise a false alarm. Since the anchored
protocol leaves them out of both training and testing, the standard pipeline
neither teaches a model to handle them nor shows whether it can.

This paper measures that omission and what it costs.

\begin{figure*}[t]
  \centering
  \includegraphics[width=\textwidth]{fig1_protocols.pdf}
  \caption{The two sampling protocols applied to one pedestrian track. (a)~The
  event-anchored protocol takes about one window per track, placed a known
  interval before a known crossing onset; the rest of the track is never
  sampled. (b)~Streaming sampling covers the whole track. Its negatives are
  dominated by windows of the same pedestrian, in the same scene, shortly before
  the onset. These are the hard temporal negatives that protocol (a) cannot
  produce. Windows are drawn separated for legibility; in practice they overlap
  heavily (length $20$, stride $3$).}
  \label{fig:protocols}
\end{figure*}

We make four contributions.
\begin{enumerate}
  \item A precise statement of both protocols, and an explanation of why their
  class ratios cannot be compared. One counts how many pedestrians eventually
  cross; the other counts how often a crossing is about to begin
  (\cref{sec:protocols}).
  \item A measurement of what the streaming negatives are made of, split by
  split. Two findings fall out of it: the PIE training and test splits differ in
  exactly the respect this paper studies, and most of the seemingly difficult
  negatives are not confusable at all (\cref{sec:omits}).
  \item A cross-protocol evaluation matrix, and a split of the
  anchored-to-streaming gap into the part that threshold re-tuning recovers and
  the part it does not. We measure the first at approximately zero
  (\cref{sec:gap}).
  \item A reformulation of the task as onset timing under right-censoring, using
  a masked discrete-time hazard objective whose horizon read-out stays
  numerically comparable with the binary baselines (\cref{sec:method}).
\end{enumerate}

We do not claim streaming evaluation as a new idea; it is well established in
the online action detection
literature~\cite{degeest2016oad,shou2018odas,gao2019startnet}. Our claim is that
the crossing-intention literature has not adopted it, that this has measurable
and substantial consequences, and that a timing formulation repairs the
underlying labeling problem rather than working around it.

% ==========================================================================
\section{Related Work}
\label{sec:related}

\subsection{Crossing Intention Estimation on PIE and JAAD}
PIE~\cite{rasouli2019pie} and JAAD~\cite{rasouli2017jaad} are the standard
datasets for this task, and the evaluation protocol released with
them~\cite{kotseruba2021benchmark} is the one most later work adopts. It fixes
per-track binary labels, observation clipped at the crossing point, and a
time-to-event interval of one to two seconds. That protocol has been inherited
largely unchanged, and reported improvements are measured against it.

Several studies revisit parts of this evaluation without changing which windows
the protocol produces. Azarmi \emph{et al.}~\cite{azarmi2024feature} divide PIE
into context sets and study which input features matter in each, finding a
strong dependence on ego-vehicle speed. Rasouli and
Kotseruba~\cite{rasouli2024diving} revise the task definitions and propose
additional classes of evaluation metric. Both work on the sample population the
anchored protocol produces; this paper questions that population itself.

The closest prior observation to ours is by Yao \emph{et
al.}~\cite{yao2021coupling}, who consider an alternative sampling setting that
keeps windows from the whole track rather than only those next to the crossing
event.
\todo[inline]{Verify against the paper body: the project notes record that this
work labels such windows by eventual crossing and treats the early ones as easy.
Only title, authors and venue have been confirmed. The contrast drawn below
depends on that detail being right.}
Our formulation differs in two ways. Windows are labeled by whether an onset
falls within a fixed horizon, not by whether the pedestrian eventually crosses;
and the early windows of eventual crossers are the main object of study rather
than samples to be set aside.
\todo[inline]{Optional before submission: add two or three recent PIE methods
for coverage. TrajFusionNet (arXiv:2508.19866) and the MDPI \emph{Algorithms}
13(12):331 pose-sequence paper are named by URL only in the project notes and
were not verified.}

\subsection{Online Action Detection and Action Start Detection}
Online action detection labels each frame of a stream using only present and
past evidence~\cite{degeest2016oad}. Its onset-focused variant, detection of
action start~\cite{shou2018odas,gao2019startnet}, must signal that an action has
begun, as early as possible, in untrimmed video where background frames
dominate. Recent work in the area is mostly
transformer-based~\cite{wang2021oadtr,xu2021lstr,zhao2022testra,wang2023mat}.
Both the task shape and the evaluation machinery built around it, including
calibrated average precision and point-level average precision with a temporal
tolerance, were designed for this regime.

That literature also describes the failure we observe. Shou \emph{et
al.}~\cite{shou2018odas} note that a start window shares its background scene
and objects with the frames before it, so its features can sit closer to the
preceding background than to the action that follows. In our setting, a window
taken just before someone steps off the curb resembles a window of the same
person waiting at the curb more than it resembles a crossing. We also note that
this literature often cites pedestrian crossing as a motivating safety-critical
application, while evaluating on television, sports and ego-vehicle maneuver
data.

\subsection{Time-to-Event Modeling}
Modeling the time until an event, when some observations end before the event is
seen, is the subject of survival
analysis~\cite{kaplan1958nonparametric,cox1972regression}. The discrete-time
form we use in \cref{sec:method}, where a hazard for each interval is trained
with a masked binary likelihood, is standard there and has been carried over to
neural estimators~\cite{lee2018deephit}.

Survival methods have already been applied to pedestrian behavior. Kalatian and
Farooq~\cite{kalatian2019deepwait} estimate how long a pedestrian waits at an
unsignalized crosswalk, using a Cox proportional hazards model whose risk
function is a neural network, trained on data from participants in a virtual
reality environment. That work shows survival modeling suits pedestrian crossing
behavior, but it addresses a different observation setting and a different
quantity: waiting time measured from trajectory-level data, rather than onset
detection from streaming in-vehicle video under the PIE protocol. A targeted
search found no prior work applying censored time-to-event modeling to
crossing-onset detection in that setting, though we do not claim the search was
exhaustive. Our contribution is therefore not the estimator. It is the
observation that streaming crossing windows already form a right-censored
time-to-event dataset, and that binary labeling discards exactly the information
that makes them one.

\subsection{Learning Under Class Imbalance}
Focal loss~\cite{lin2017focal} and class-balanced
re-weighting~\cite{cui2019classbalanced} are the standard responses to a rare
positive class. \cref{sec:gap} gives measured reasons to expect them to fall
short here: they mainly move the decision boundary, whereas what degrades under
protocol shift is the ranking itself. We therefore treat them as comparison
baselines rather than as the proposed remedy.

% ==========================================================================
\section{Problem Formulation}
\label{sec:protocols}

\subsection{Event-Anchored Sampling}
Let a pedestrian track consist of annotated frames $1,\dots,n$. If the
pedestrian crosses, let $\tau$ be the frame at which the crossing begins. The
anchored protocol stops observation at $\tau$, or at the last visible frame for
non-crossers, and emits a small number of observation windows of length $T$
whose final frame falls in a fixed time-to-event interval before $\tau$, usually
$30$ to $60$ frames. The label belongs to the track: this pedestrian either
crosses or does not. The positive share therefore reflects how many pedestrians
in the dataset eventually cross, which for PIE gives roughly $2.5{:}1$.

\subsection{Streaming Sampling}
The streaming protocol places a window at every stride-$s$ position along the
whole track. For a window whose last observed frame is $t$, the label is
\begin{equation}
  y_t = \mathbf{1}\left[\, \exists\ \text{onset in } (t,\, t+H] \,\right],
  \label{eq:binlabel}
\end{equation}
where $H$ is the anticipation horizon. Windows in which a crossing is already
underway during observation are discarded, since a detector cannot anticipate an
onset that has already happened. The positive share now reflects how often, per
observed moment, a crossing is imminent. On PIE this is $2.9\%$ of windows, a
ratio of $33.9{:}1$.

\subsection{Why the Two Ratios Are Not Comparable}
The two ratios measure different things. The anchored $2.5{:}1$ is a balance
across tracks: it says how many pedestrians eventually cross, because the
protocol draws about one window from each. The streaming $33.9{:}1$ is a rate
over time: it says how often an imminent crossing occurs per observed moment.
They are not two estimates of one number, and moving between them is not a
difficulty dial. The gap between them follows from which windows each protocol
creates in the first place.

One reasonable objection is that streaming sampling is just anchored sampling
with a long time-to-event. It is not. A long-time-to-event anchored sample is
still one window per track, placed a known interval from a known event.
Streaming sampling uses no anchor and no knowledge of the event, and its
negatives are dominated by a category the anchored protocol does not generate.
\cref{sec:omits} provides the measurements that show this.

\subsection{Four Kinds of Streaming Negative}
\label{sec:taxonomy}
Streaming negatives fall into four kinds, shown in \cref{fig:protocols}(b).
\begin{itemize}
  \item \textbf{Genuine non-crossers.} People waiting, or walking along the
  sidewalk, who do not cross. Present under both protocols.
  \item \textbf{Hard temporal negatives.} People who will cross, but not within
  $H$. The person, clothing, street and camera are the same as in the positive
  windows that follow; only the position of the cutoff makes the label differ.
  \item \textbf{Already crossed.} Windows where the crossing happened before the
  observation began. These pedestrians are not at risk of a first crossing at
  all, yet they receive a negative label and sit alongside genuine
  non-crossers.
  \item \textbf{Uninformative windows.} Occluded, distant or otherwise
  low-information observations.
\end{itemize}
The second and third kinds are the ones the anchored protocol structurally
omits. The third is not merely a hard case: it is a labeling error the binary
formulation gives no way to avoid.

% ==========================================================================
\section{What the Streaming Window Population Contains}
\label{sec:omits}

\subsection{Setup}
We use PIE with the standard split assignment: sets 01, 02 and 06 for training,
sets 04 and 05 for validation, and set 03 for testing. Streaming windows are
$T = 20$ frames at stride $3$, with horizon $H = 32$ frames, about $1.07$\,s at
$30$\,fps. Windows whose future is not fully observed are discarded rather than
labeled negative, since we cannot establish that no crossing occurred in an
interval nobody watched. \cref{tab:stats} gives the resulting population. Every
measurement in this section comes straight from the PIE annotation files and can
be reproduced without image data.

\begin{table}[t]
  \centering
  \caption{Streaming window population on PIE.}
  \label{tab:stats}
  \begin{tabular}{l S[table-format=5.0] S[table-format=2.1] S[table-format=2.1] S[table-format=1.1]}
    \toprule
    Split & {$N$} & {actions=1 (\%)} & {looks=1 (\%)} & {crosses=1 (\%)} \\
    \midrule
    train & 88214 & 43.4 & 10.5 & 2.9 \\
    val   & 20490 & 40.3 &  6.7 & 2.8 \\
    test  & 69875 & 42.2 &  9.8 & 3.1 \\
    \bottomrule
  \end{tabular}
\end{table}

\subsection{What the Negatives Are Made Of}
Matching each window against the crossing annotation of its own track gives the
breakdown in \cref{fig:composition} and \cref{tab:composition}. About one
training window in four, and two test windows in five, show a pedestrian who
does cross, only later than \cref{eq:binlabel} asks about. Such windows
outnumber the true positives by $8.9{:}1$ in training and $12.8{:}1$ in test.

\begin{figure}[t]
  \centering
  \includegraphics[width=\columnwidth]{fig2_composition.pdf}
  \caption{What streaming windows are made of, by split. The hard temporal
  negatives, shown as ``will cross, later'', are the population the anchored
  protocol does not generate. Training and test differ sharply here: $25.7\%$
  against $39.2\%$.}
  \label{fig:composition}
\end{figure}

\begin{table}[t]
  \centering
  \caption{Streaming negative composition, as a share of each split (\%).}
  \label{tab:composition}
  \begin{tabular}{l S[table-format=1.1] S[table-format=2.1] S[table-format=2.1] S[table-format=2.1]}
    \toprule
    Split & {positive} & {will cross, later} & {already crossed} & {never crosses} \\
    \midrule
    train & 2.9 & 25.7 & 7.2 & 64.3 \\
    val   & 2.8 & 15.2 & 4.6 & 77.5 \\
    test  & 3.1 & 39.2 & 4.3 & 53.4 \\
    \bottomrule
  \end{tabular}
\end{table}

Two things in this table deserve attention.

First, the splits are not built the same way. Test holds about half again as
many hard temporal negatives as training, $39.2\%$ against $25.7\%$, and
correspondingly fewer easy ones. A drop from training to test performance on
this data therefore cannot be read as a generalization gap alone, because the
test split is harder in exactly the respect under study. Validation is the
easiest of the three at $15.2\%$, which means thresholds chosen on validation
are optimistic for test in a way that a correctly run no-peeking protocol does
not fix.

Second, most of the seemingly difficult mass is not confusable. Grouping the
hard temporal negatives by how far the crossing still is
(\cref{fig:onsethist}), $68.7\%$ lie more than five seconds ahead. Someone
walking down a street who crosses several seconds later is not a confusing
example in any useful sense. The genuinely confusable group is the windows whose
onset falls in the two seconds just past the horizon, between $1.1$\,s and
$3.2$\,s: $4{,}186$ training windows, $1.65$ times the size of the positive set.
That is the population a targeted method has to work on, and it is large enough
to train against without swamping everything else.

\begin{figure}[t]
  \centering
  \includegraphics[width=\columnwidth]{fig3_time_to_onset.pdf}
  \caption{Hard temporal negatives in the training split, grouped by how long
  remains until the crossing starts. Two thirds lie more than five seconds ahead
  and are not confusable. The confusable band, the two shaded groups, is about
  the size of the positive set.}
  \label{fig:onsethist}
\end{figure}

\subsection{Why Not Simply Predict Further Ahead}
\label{sec:horizon}
The obvious answer to a $33.9{:}1$ imbalance is to enlarge $H$. Relabeling the
same windows at a range of horizons (\cref{tab:horizon}) shows this improves the
nominal ratio considerably: the imbalance falls to $5.8{:}1$ at five seconds,
and the confusable band shrinks from $1.65$ to $0.26$ times the positive set. We
keep $H = 32$ nonetheless, for three reasons.

\begin{table}[t]
  \centering
  \caption{Relabeling the same windows at longer horizons (training split). The
  confusable band is the set of windows whose onset falls in the two seconds
  after the horizon.}
  \label{tab:horizon}
  \begin{tabular}{l S[table-format=5.0] S[table-format=5.0] l S[table-format=4.0] S[table-format=2.0]}
    \toprule
    Horizon & {usable} & {positives} & imbalance & {band} & {lost (\%)} \\
    \midrule
    1.1\,s\ \emph{(ours)} & 88214 &  2530 & 33.9:1 & 4186 &  0 \\
    1.5\,s                & 85297 &  3471 & 23.6:1 & 4001 &  3 \\
    2.0\,s                & 82020 &  4481 & 17.3:1 & 3816 &  7 \\
    3.0\,s                & 75767 &  6367 & 10.9:1 & 3435 & 14 \\
    5.0\,s                & 65326 &  9604 &  5.8:1 & 2499 & 26 \\
    10.0\,s               & 49325 & 14599 &  2.4:1 & 1496 & 44 \\
    \bottomrule
  \end{tabular}
\end{table}

\textbf{A longer horizon deletes the windows that matter most.} It needs more
observed future, so $26\%$ of windows fall away at five seconds and $44\%$ at
ten. They fall away by track length, so short tracks go first, and short tracks
are pedestrians who appear abruptly or are quickly occluded. Those are precisely
the cases a short-horizon detector exists to cover. The improved balance is
bought by deleting the hardest part of the dataset.

\textbf{The confusion moves rather than disappears.} The ratio improves because
windows that were hard negatives at one second become positives at five. The
model still cannot tell a window $4.9$\,s before onset from one $5.1$\,s before
onset. It now has to answer yes on the first, on evidence that is not in the
frame. The errors turn from false alarms into misses.

\textbf{The horizon decides which phenomenon is being detected.} At about one
second, what is visible is the body starting to move: weight shifting, a foot
lifting, the torso turning. At five seconds no such cue exists yet, and what is
visible instead is where the person is heading relative to the road. These are
two different problems needing two different feature sets, so a longer horizon
replaces the task rather than extending it.

We note that a ten-second horizon gives $2.4{:}1$, close to the anchored
protocol's $2.5{:}1$. The two arise by different mechanisms and the match is
coincidental, but it does show how far the anchored ratio reflects a sampling
choice rather than a property of the task.

% ==========================================================================
\section{Cross-Protocol Evaluation}
\label{sec:gap}

\subsection{Model and Training Setup}
All runs use the same multimodal architecture: a hierarchical windowed-attention
vision transformer over context crops, a pose and kinematics stream in which
$23$ whole-body keypoints~\cite{jin2020coco} are reduced to a $58$-dimensional
per-frame vector, and a cross-attention module in which the kinematic stream
queries the visual features. Per-task classification heads run at a shared model
dimension of $128$. The visual backbone is a frozen ImageNet-pretrained
TinyViT-5M~\cite{wu2022tinyvit}. All four runs share a learning rate of
$10^{-4}$, weight decay of $10^{-5}$, an effective batch size of $32$ (batch
size $4$ with $8$ accumulation steps), a warmup--cosine learning-rate schedule
with four warmup epochs, and seed $42$.

\subsection{The Cross-Protocol Matrix}
Each model is trained under one protocol and evaluated under both. Thresholds
are chosen on validation and applied unchanged to test; same-split threshold
sweeps are computed for diagnosis but never reported as results.
\cref{tab:matrix} gives the matrix for two models, one trained on all three
annotation tasks and one trained on the crossing task alone.

\begin{table}[t]
  \centering
  \caption{Cross-protocol results for the crossing task. Cells report
  AUC $\cdot$ raw F1 $\rightarrow$ threshold-tuned F1.}
  \label{tab:matrix}
  \setlength{\tabcolsep}{3pt}
  \scriptsize
  \begin{tabular}{ll cc}
    \toprule
    & & \multicolumn{2}{c}{test protocol} \\
    \cmidrule(lr){3-4}
    Model & train protocol & anchored & streaming \\
    \midrule
    \multirow{2}{*}{A (3-task)}
      & anchored  & 0.873 $\cdot$ 0.689 $\rightarrow$ 0.694 & \textbf{0.540 $\cdot$ 0.064 $\rightarrow$ 0.064} \\
      & streaming & 0.514 $\cdot$ 0.186 $\rightarrow$ 0.322 & 0.742 $\cdot$ 0.211 $\rightarrow$ 0.190 \\
    \midrule
    \multirow{2}{*}{B (crossing-only)}
      & anchored  & 0.880 $\cdot$ 0.716 $\rightarrow$ 0.721 & \textbf{0.529 $\cdot$ 0.064 $\rightarrow$ 0.062} \\
      & streaming & 0.679 $\cdot$ 0.258 $\rightarrow$ 0.428 & 0.784 $\cdot$ 0.249 $\rightarrow$ 0.225 \\
    \bottomrule
  \end{tabular}
\end{table}

Scores within a single protocol are in line with published results, and the
off-diagonal entries carry the finding. A model trained under the anchored
protocol reaches $0.873$ AUC on the anchored test set and $0.540$ on the
streaming test set, which is close to chance. The same collapse appears in the
crossing-only model, $0.880$ down to $0.529$, so interference between the
multi-task heads is not the cause.

\subsection{Splitting the Deployment Gap}
The deployment gap can be split into the part that re-tuning the threshold
recovers and the part it does not. With every quantity measured in tuned F1 on
the streaming test column:
\begin{align}
  \Gtot   &= \mathrm{F1}^{\text{tuned}}_{\text{str}\to\text{str}}
             - \mathrm{F1}^{\text{raw}}_{\text{anc}\to\text{str}}, \label{eq:gtot}\\
  \Gprior &= \mathrm{F1}^{\text{tuned}}_{\text{anc}\to\text{str}}
             - \mathrm{F1}^{\text{raw}}_{\text{anc}\to\text{str}}, \label{eq:gprior}\\
  \Ghard  &= \mathrm{F1}^{\text{tuned}}_{\text{str}\to\text{str}}
             - \mathrm{F1}^{\text{tuned}}_{\text{anc}\to\text{str}}, \label{eq:ghard}
\end{align}
so that $\Gtot = \Gprior + \Ghard$ by construction. Here $\Gprior$ is what an
engineer gains simply by re-choosing the operating point at the streaming base
rate, and $\Ghard$ is what remains, which we attribute to an ability to separate
hard negatives that anchored training never builds.

\begin{table}[t]
  \centering
  \caption{Splitting the anchored-to-streaming gap.}
  \label{tab:decomp}
  \begin{tabular}{l S[table-format=+1.3] S[table-format=+1.3]}
    \toprule
    & {Model A} & {Model B} \\
    \midrule
    $\Gtot$   & +0.126 & +0.161 \\
    $\Gprior$ & -0.001 & -0.001 \\
    $\Ghard$  & +0.126 & +0.162 \\
    \bottomrule
  \end{tabular}
\end{table}

\cref{tab:decomp} gives the result. For both models $\Gprior \approx 0$ and
$\Ghard \approx \Gtot$. Essentially none of the deployment gap is recoverable by
recalibration, and essentially all of it is a failure to separate hard
negatives.

The AUC figures explain why. Re-tuning a threshold slides a decision boundary
along an existing ranking. When the ranking itself has fallen to $0.540$, barely
better than ordering the pedestrians at random, there is no threshold left to
find. The argument reaches past thresholds: any method whose main effect is to
move the operating point, including class re-weighting and focal
loss~\cite{lin2017focal,cui2019classbalanced}, runs into the same limit.

\subsection{Threats to Validity}
\label{sec:threats}
Four qualifications travel with \cref{tab:matrix} and should stay attached
wherever these numbers are quoted.

\textbf{Training set size is confounded with protocol.} The anchored training
set holds about $4.9$k windows and the streaming set about $88$k, roughly
$18$ times more. $\Ghard$ therefore mixes the effect of the training protocol
with the effect of training set size, so results should be described as
consistent with a hard-negative gap rather than as establishing one.
\todo[inline]{Highest-priority outstanding experiment: a streaming run
subsampled to about $4.9$k windows. One training run, and it would let the
central claim be stated as demonstrated rather than as consistent with the
evidence.}

\textbf{F1 is unstable at this base rate.} With one positive in $34$, a handful
of samples crossing the threshold moves F1 noticeably. The AUC drop, $0.88$ to
$0.53$ in both models, is the more reliable measurement, and average precision
and precision at fixed recall are the appropriate replacements.

\textbf{Model B has a configuration mismatch.} Its two halves were trained with
different sampler settings: the streaming half oversampled the rare crossing
windows and the anchored half did not. We therefore cite Model B only for the
point that multi-task interference is not the cause of the gap, never for the
size of $\Ghard$.

\textbf{Every result is from a single seed.} Each entry is one run at seed
$42$, with no variance estimate.
\todo[inline]{The multi-seed protocol (three seeds, mean and standard deviation)
has not been applied to any entry in \cref{tab:matrix}.}

% ==========================================================================
\section{Onset Timing Under Right-Censoring}
\label{sec:method}

\subsection{Where the Binary Label Goes Wrong}
\cref{eq:binlabel} gets two situations wrong by construction.

A window whose crossing begins at $H + 5$ frames is labeled negative and grouped
with pedestrians who never cross, even though it looks almost identical to a
window at $H - 5$ frames, which is labeled positive. This is less a hard
learning problem than an ill-posed one, and it comes from where we placed the
cutoff rather than from the data.

A window whose footage ends before $H$ cannot honestly be given either label.
Calling it negative asserts an observation nobody made; discarding it, as we do
in \cref{sec:omits}, throws away information that is genuinely there. Neither is
right, because the true statement, that no crossing happened before the footage
ran out, cannot be written in a single bit.

Both situations are what a right-censored time-to-event dataset looks like when
it is forced into binary classification.

\subsection{Three Numbers Kept per Window}
For each streaming window we keep three quantities alongside the binary label:
\begin{itemize}
  \item $o$, how many frames pass between the last observed frame and the next
  crossing onset, or $-1$ if no onset occurs in the remaining footage;
  \item $f$, how many frames of future footage exist after the window;
  \item $c$, whether the track contains a crossing at any point.
\end{itemize}
In survival-analysis terms $o$ is the event time, $f$ is the censoring time, and
$c$ tells apart the two situations in which no crossing is seen. The
correspondence is exact, not an analogy.

\subsection{Discrete-Time Hazard}
Divide the next $L$ frames into $K = L/w$ bins of width $w$. Let $h_k$ be the
hazard of bin $k$: the probability that a first crossing begins in that bin,
given that none has begun earlier. The model produces $K$ logits from the pooled
representation. The probability of an onset within a horizon spanning
$K_H = H/w$ bins then follows from the survival product:
\begin{equation}
  P(\text{onset} \le H) = 1 - \prod_{k=0}^{K_H - 1}\left(1 - h_k\right).
  \label{eq:readout}
\end{equation}
\cref{eq:readout} is what keeps the comparison with existing results alive,
because it answers exactly the question \cref{eq:binlabel} asks. The timing
model and the binary baselines report the same quantity under the same metric.

\subsection{Supervision When the Future Is Only Partly Seen}
The three numbers determine a per-bin target $y_k$ and a per-bin mask $m_k$, as
shown in \cref{fig:cases}.
\begin{enumerate}
  \item \textbf{Onset within range} ($0 \le o < L$). With $e = \lfloor o/w
  \rfloor$, bins $0$ to $e-1$ are supervised as $0$ and bin $e$ as $1$. Every
  bin after $e$ is masked, because once a first crossing has begun, later bins
  say nothing about a first crossing.
  \item \textbf{Onset beyond range} ($o \ge L$). All $K$ bins are supervised as
  $0$. This is honest: the crossing was seen, and it did not fall inside the
  window's look-ahead.
  \item \textbf{Censored} ($o < 0$, with some future seen). Bins covered by $f$
  are supervised as $0$ and the rest are masked. This is the case the binary
  label cannot express.
  \item \textbf{Already crossed} ($o < 0$ and $c$). The window is dropped, since
  this pedestrian is not at risk of a first crossing. This is the case the
  binary label cannot get right.
\end{enumerate}

\begin{figure}[t]
  \centering
  \includegraphics[width=\columnwidth]{fig4_hazard_cases.pdf}
  \caption{The four supervision cases. Shaded bins contribute to the loss and
  dashed bins are masked, receiving no gradient; case~4 leaves the risk set
  altogether. Case~3 is the situation the binary label cannot express, and
  case~4 is the one it cannot get right.}
  \label{fig:cases}
\end{figure}

\subsection{Training Objective}
The loss is a masked per-bin binary cross-entropy:
\begin{equation}
  \mathcal{L}_{\text{onset}} = -\sum_{k=0}^{K-1} m_k
    \Big[\, y_k \log h_k + (1 - y_k)\log(1 - h_k) \,\Big].
  \label{eq:loss}
\end{equation}
The masking is the substantive part. A bin nobody observed must receive exactly
zero gradient, not a small gradient toward zero, and we check this with
automatic differentiation rather than by reading the code.

Two constraints are enforced in the implementation rather than left as advice.
First, $L > H$ strictly. A look-ahead no wider than the reported horizon would
still hand a flat negative label to a crossing at $H + 5$, which is the very
defect the formulation exists to remove. Second, the read-out horizon must match
the horizon used by the binary baselines, or \cref{eq:readout} quietly answers a
different question under the same metric name. A test enforces the second by
reproducing the data generator's own binary label from the timing read-out over
a full track.

\subsection{Choosing the Look-Ahead and Bin Width}
The reported horizon stays at $H = 32$ frames, for the reasons in
\cref{sec:horizon}. The look-ahead $L$ is a separate choice, and we size it
against the composition measured in \cref{sec:omits}. Setting $L = 96$ frames,
or $3.2$\,s, covers about $92\%$ of the confusable band found there, against
$47\%$ at $L = 60$. Windows whose onset falls inside $L$ carry a real ``crossing
begins here'' signal instead of a flat negative; at $L = 96$ that applies to
roughly $6{,}400$ training windows, against $2{,}530$ at $L = H$. That is close
to a threefold increase in the windows carrying any positive supervision, in a
task with a $2.9\%$ positive rate. Going wider buys little, since two thirds of
the hard negative mass sits more than five seconds ahead.

A bin width of $w = 4$ gives $K = 24$ bins, keeps $8$ bins inside the reported
horizon, makes positives about four times denser per bin, and costs $133$\,ms of
timing resolution. Both settings affect training only; neither changes what is
reported.

\subsection{Three Configurations of One Objective}
The formulation supports three configurations, chosen by loss weights alone with
no change to the code. In the \emph{auxiliary} configuration the binary head
still produces the reported output and the hazard term acts as a regularizer. In
the \emph{pure} configuration the hazard is the only objective and
\cref{eq:readout} produces the reported output. The \emph{hedged} configuration
adds a direct read-out term. The first two are enough to test the methodological
claim; the third exists because in the second, the quantity being optimized is
not exactly the quantity being reported.

One scaling point matters in practice. \cref{eq:loss} sums over the bins a
window actually observed, so a fully observed window contributes about $17$ nats
at initialization, against $5.5$ for a window carrying only the $32$ frames of
future that generation guarantees. The term therefore starts far larger than a
per-task cross-entropy and has to be turned down when it runs alongside the
binary objective.

\subsection{The Failure Mode to Watch For}
With $K$ bins, each bin sees a positive only about $2.9\%/K$ of the time. The
head can therefore drive the loss down by predicting $h \approx 0$ everywhere,
which makes \cref{eq:readout} collapse toward zero for every window. Watching
the mean hazard will not catch this, because a low mean hazard is correct. What
to watch instead is the spread of the read-out across validation windows: a
narrow spike at the base rate means the head has learned the average and nothing
more. Setting $w = 4$ is the mitigation already applied; if an early run still
shows this behavior, $w = 8$ halves $K$ again.

% ==========================================================================
\section{Experiments}
\label{sec:experiments}

\subsection{Evaluation Protocol}
Thresholds are chosen on the validation split and applied unchanged to test. We
report AUC as the measure of ranking ability and F1 at the tuned threshold as
the operating-point measure, and treat AUC as primary for the reasons in
\cref{sec:threats}.
\todo[inline]{Adopt the online action detection metric suite before the method
results arrive: calibrated average precision, and point-level average precision
with a temporal tolerance (\cref{sec:related}). F1 at $34{:}1$ is too unstable
to tell configurations apart, and using the neighboring field's metrics is what
makes these numbers comparable with theirs.}

\subsection{Baselines}
The four runs in \cref{tab:matrix} are the baselines. They share the
architecture and optimizer settings of \cref{sec:gap}, differ only in training
protocol and task set, and were evaluated under both protocols.

\subsection{Onset Timing Results}
\begin{table}[t]
  \centering
  \caption{Onset-timing configurations on the streaming test split. The read-out
  is \cref{eq:readout} evaluated at $H = 32$, and is directly comparable with
  the baseline row.}
  \label{tab:onset}
  \begin{tabular}{l cccc}
    \toprule
    Configuration & AUC & AP & tuned F1 & read-out spread \\
    \midrule
    binary baseline (str$\to$str) & 0.742 & --- & 0.190 & n/a \\
    auxiliary hazard              & ---   & --- & ---   & --- \\
    pure hazard                   & ---   & --- & ---   & --- \\
    hedged                        & ---   & --- & ---   & --- \\
    \bottomrule
  \end{tabular}
\end{table}
\todo[inline]{\cref{tab:onset} is empty. The onset head, loss, targets and
read-out are implemented and unit-tested, including the zero-gradient property
and the match between the read-out and the stored label, but nothing has been
trained. Required order: (i)~metadata backfill on the existing databases,
(ii)~a short run to confirm the read-out is not degenerate, (iii)~the auxiliary
configuration, (iv)~the pure configuration. The AP column needs the metric suite
noted above.}

\subsection{Planned Ablations}
Besides the three configurations, we plan to vary the look-ahead over
$L \in \{60, 96, 150\}$ and the bin width over $w \in \{4, 8\}$ to check the
sizing argument of \cref{sec:method}; to run focal and class-balanced losses on
the binary objective as representatives of the move-the-boundary family; and to
compare frame-pooling rules, which are configurable and so cost nothing to
implement.

% ==========================================================================
\section{Discussion}
\label{sec:discussion}

\subsection{How Far Ahead Crossing Is Predictable at All}
Consider a pedestrian walking along the sidewalk exactly like everyone else, who
then turns abruptly and steps into the road. Three seconds earlier, the
information is not in the video, either because the decision has not been made
or because it has not yet reached the body. No feature, architecture or loss
function can recover it.

Some unknown share of the confusable band is therefore not hard but impossible,
and a claim to handle hard negatives in general would overstate what is
achievable. We accept three consequences. The goal is not accuracy on that band,
but confidence where evidence exists and uncertainty where it does not, which is
a matter of ranking and calibration rather than classification; this also
matches the failure we measured, where the ordering degraded rather than the
operating point. The binary label is a modeling choice rather than a property of
the task, and \cref{sec:method} removes it. Finally, where this limit sits can
itself be measured: given a trained model, separability can be evaluated against
time to onset, and the point at which it reaches chance can be reported.
\todo[inline]{This predictability-limit measurement needs only one trained model
and a stratified evaluation. No prior work in the area appears to report it, and
it turns the strongest reviewer objection into a reported result. Worth
including before submission.}

\subsection{Limitations}
Beyond the four qualifications in \cref{sec:threats}, all results come from PIE,
recorded in one city in daylight and clear weather; replication on JAAD would be
needed to show the effect is not specific to this dataset. The streaming
training runs are also unstable, with validation loss spiking and recall
snapping to $1.0$ on several epochs, meaning the model briefly collapses into
calling everything a crossing. The streaming-trained entries therefore rest on a
selected best epoch that may not be a settled optimum. Whether that instability
is an artifact to fix or a finding about training at realistic base rates is
still open.

\subsection{What This Means for Evaluation Practice}
If $\Ghard \approx \Gtot$ holds more widely, then a model's standing under the
anchored protocol says little about how it will behave in a stream, and the
standard protocol has been systematically under-training and under-testing the
safety-critical case. That is a claim about evaluation practice, and it is cheap
to check, because the streaming sampler is a relabeling of annotations that
already exist rather than a new data collection. Recent work revisiting the
definitions and metrics used on PIE and JAAD~\cite{rasouli2024diving} shows that
the adequacy of the standard evaluation is already an open concern; our results
suggest the sampling protocol, not only the metrics, deserves part of that
attention.

% ==========================================================================
\section{Conclusion}
\label{sec:conclusion}

The event-anchored protocol and the streaming condition are not two difficulty
settings of one task, because they contain different negatives. We measured what
the anchored protocol leaves out, finding that a quarter of streaming training
windows and two fifths of test windows show pedestrians who cross later, and we
showed that a model trained under the anchored protocol loses almost all ranking
ability on the streaming test set, with threshold re-tuning recovering
essentially none of the loss. Since part of the difficulty lies in the labeling
rather than in the model alone, we proposed treating the task as onset timing
under right-censoring. A window whose future footage runs out is then censored
rather than negative, a pedestrian who has already crossed leaves the risk set
instead of counting as a non-crosser, and the reported horizon probability stays
directly comparable with the binary baselines.

\todo[inline]{Pre-submission checklist: (1)~fill \cref{tab:onset};
(2)~matched-size streaming control; (3)~verify the Yao \emph{et al.} sampling
description against the paper body; (4)~multi-seed the cross-protocol matrix;
(5)~predictability-limit measurement; (6)~adopt the OAD metric suite.}

\bibliographystyle{IEEEtran}
\bibliography{refs}

\end{document}
